\documentclass{article}

\usepackage{float}
\restylefloat{table}

\usepackage{booktabs}

\title{Team Productivity: POC\\\progname}

\author{\authname}

\date{}

\input{../Comments.text}
\input{../Common.text}

\begin{document}

\maketitle

This document summarizes the contributions of each team member up to the POC
Demo.  The time period of interest is the time between the beginning of the term
and the POC demo.

\wss{Please do not delete sections or modify the table formats.  Standardization
helps the instructors and TAs to review these documents. You may add sections to
the report.}

\wss{Remove the instructor guidelines/comments in the final version.}

\wss{Fill in the requested information, like project url, TA name, supervisor name, etc.}

\section{Demo Plans}

The main feasibility risk is whether Unity can render a first-person view in
a Meta Quest headset while receiving controller inputs that drive the simulated
drone. The prototype initially appeared only on the monitor and received no
Quest controller input. We enabled Unity's OpenXR loader and Quest controller
profiles, then verified headset rendering, head tracking, and thumbstick and
trigger input over USB Quest Link.
\href{https://github.com/SoftwareMandems/FPVR-Game/pull/12}{The Quest Link
integration pull request} records these results.

During the demo, we will launch the Unity training scene, show how the drone
responds to Quest controller input on the gate course, and reset the run.
We will show desktop keyboard or gamepad controls as a fallback for headset or
controller integration problems. Using supported controllers lets us demonstrate
basic flight while a custom radio link remains under investigation. We started
with a configurable flight model so we can tune the drone's response, and the
prototype lets us disable costly rendering effects if frame rate drops. The
remaining work includes correcting the
\href{https://github.com/SoftwareMandems/FPVR-Game/issues/13}{known control
mapping issue} and measuring frame rate and input latency on the target
hardware.

\section{Team Meeting Attendance}

\wss{For each team member how many team meetings have they attended over the
time period of interest.  This number should be determined from the meeting
issues in the team's repo.  The first entry in the table should be the total
number of team meetings held by the team.}

\begin{table}[H]
\centering
\begin{tabular}{ll}
\toprule
\textbf{Student} & \textbf{Meetings}\\
\midrule
Total & 3\\
Ibrahim Tariq & 3\\
John Cho & 3\\
Sankalpa Dhakal & 3\\
Vikram Chandar & 3\\
Arjun Sasidaran & 3\\
\bottomrule
\end{tabular}
\end{table}

\wss{If needed, an explanation for the counts can be provided here.}

As of November 15, 2026, we held three team meetings, on October 8,
October 19, and November 2. Since October 19, team meetings have followed an
every-other-Monday schedule, from 3:30 pm to 4:30 pm. All five team members
attended every meeting. The closed meeting records are
\href{https://github.com/SoftwareMandems/FPVR/issues/75}{\#75},
\href{https://github.com/SoftwareMandems/FPVR/issues/283}{\#283}, and
\href{https://github.com/SoftwareMandems/FPVR/issues/284}{\#284}.

\section{Supervisor/Stakeholder Meeting Attendance}

\wss{For each team member how many supervisor/stakeholder team meetings have
they attended over the time period of interest.  This number should be determined
from the supervisor meeting issues in the team's repo.  The first entry in the
table should be the total number of supervisor and team meetings held by the
team.  If there is no supervisor, there will usually be meetings with
stakeholders (potential users) that can serve a similar purpose.}

\noindent \textbf{Supervisor's Name: } Dr. Douglas Down

\begin{table}[H]
\centering
\begin{tabular}{ll}
\toprule
\textbf{Student} & \textbf{Meetings}\\
\midrule
Total & 2\\
Ibrahim Tariq & 2\\
John Cho & 2\\
Sankalpa Dhakal & 2\\
Vikram Chandar & 2\\
Arjun Sasidaran & 2\\
\bottomrule
\end{tabular}
\end{table}

\wss{If needed, an explanation for the counts can be provided here.}

Supervisor meetings are scheduled every other Wednesday beginning October 7,
2026. As of November 15, two meetings were held, and all five team members
attended both. The third planned meeting was rescheduled because of a lack of
updates and scheduling conflicts. The two completed meetings are recorded in
\href{https://github.com/SoftwareMandems/FPVR/issues/73}{\#73} and
\href{https://github.com/SoftwareMandems/FPVR/issues/151}{\#151}.

\section{Lecture Attendance}

\wss{For each team member how many lectures have they attended over the time
period of interest.  This number should be determined from the lecture issues in
the team's repo. You can find the number of lectures in the time period of
interest by looking at the
\href{https://calendar.google.com/calendar/u/0/embed?src=rnboqiaki1k2la7rpu3bn0um58@group.calendar.google.com&ctz=America/Toronto}
{Google calendar} for the capstone course. The first entry in the
table should be the total number of lectures held for the class during the period of interest.}

\wss{NOTE: There will be approximately 13 lectures between the start of class
and the POC demos}

\begin{table}[H]
\centering
\begin{tabular}{ll}
\toprule
\textbf{Student} & \textbf{Lectures}\\
\midrule
Total & 11\\
Ibrahim Tariq & 8\\
John Cho & 10\\
Sankalpa Dhakal & 9\\
Vikram Chandar & 9\\
Arjun Sasidaran & 8\\
\bottomrule
\end{tabular}
\end{table}

\wss{If needed, an explanation for the lecture attendance can be provided here.}

Eleven lectures had been held as of November 15, 2026. Although individual
attendance varied, every lecture was attended by several team members so that the
group received the relevant course information.

\section{TA Document Discussion Attendance}

\wss{For each team member how many of the informal document discussion meetings
with the TA were attended over the time period of interest. The first row is the
total number of discussion meetings scheduled by the TA with the team.}

\noindent \textbf{TA's Name:} Rebecca Di Filippo

\begin{table}[H]
\centering
\begin{tabular}{ll}
\toprule
\textbf{Student} & \textbf{Meetings}\\
\midrule
Total & 4\\
Ibrahim Tariq & 4\\
John Cho & 4\\
Sankalpa Dhakal & 4\\
Vikram Chandar & 4\\
Arjun Sasidaran & 4\\
\bottomrule
\end{tabular}
\end{table}

\wss{If needed, an explanation for the attendance can be provided here.}

Four TA document discussions had been held as of November 15, 2026. All five
team members attended all four discussions.

\section{Commits, Lines Added and Lines Deleted, Pull Requests}

\begin{table}[H]
\centering
\begin{tabular}{lllllll}
\toprule
\textbf{Student} & \textbf{Commits} & \textbf{Percent} & \textbf{Lines Added} &
\textbf{Lines Deleted} & \textbf{Pull Requests}\\
\midrule
Total & 427 & 100\% & 30,273 & 7,412 & 119\\
Ibrahim Tariq & 93 & 21.78\% & 13,978 & 1,276 & 22\\
John Cho & 81 & 18.97\% & 1,302 & 346 & 25\\
Sankalpa Dhakal & 88 & 20.61\% & 4,219 & 2,518 & 23\\
Vikram Chandar & 83 & 19.44\% & 3,555 & 930 & 24\\
Arjun Sasidaran & 82 & 19.20\% & 7,219 & 2,342 & 25\\
\bottomrule
\end{tabular}
\end{table}

We used \href{https://github.com/SoftwareMandems/FPVR/graphs/contributors}{GitHub
Insights} to review repository contributions and verified the figures against
GitHub pull request records and the main-branch snapshot
\href{https://github.com/SoftwareMandems/FPVR/commit/b9ffc85a6e601fa25f906428faeb0058c3502530}{\texttt{b9ffc85}}.
These figures cover September 17 through November 19, 2026. Each team-authored
commit reachable from \texttt{main}, including merge commits, counts once
under its primary author. Git author email determines attribution. Seven
commits display the name \texttt{IbrahimT} but use Sankalpa's GitHub email
(five) or Vikram's (two). The team total excludes 80 automated GitHub Actions
commits. The coauthored reflection commits
\href{https://github.com/SoftwareMandems/FPVR/commit/b18fd072e1377cc2e5fa4c36acc694fb758ccdb6}{\texttt{b18fd07}}
and
\href{https://github.com/SoftwareMandems/FPVR/commit/31efc15f7c197ef68895da722466d494bc65d86d}{\texttt{31efc15}}
are counted under Arjun as primary author.

Line totals count text additions and deletions from non-merge commits.
Ibrahim's
\href{https://github.com/SoftwareMandems/FPVR/commit/2dfed2dc67f478f6178f4257a1a28202987b39cb}{initial repository import}
accounts for 11,354 of his 13,978 added lines (81.23\%), mostly course templates
and repository support files. He added 2,624 lines in subsequent commits.

The 119 team-authored pull requests include 111 merged, four open, and four
closed without merging. Ibrahim authored one of the open pull requests and
Arjun authored three. Commits found only on unmerged branches are outside the
\texttt{main} figures. The pull request count includes the six earlier
productivity report pull requests.

\section{Issue Tracker}

\begin{table}[H]
\centering
\begin{tabular}{lll}
\toprule
\textbf{Student} & \textbf{Authored (O+C)} & \textbf{Assigned (C only)}\\
\midrule
Ibrahim Tariq & 55 & 26 \\
John Cho & 52 & 26 \\
Sankalpa Dhakal & 12 & 28 \\
Vikram Chandar & 11 & 31 \\
Arjun Sasidaran & 44 & 22 \\
\bottomrule
\end{tabular}
\end{table}

We used \href{https://github.com/SoftwareMandems/FPVR/pulse}{GitHub Insights
Pulse} to review issue activity and the
\href{https://github.com/SoftwareMandems/FPVR/issues}{GitHub issue tracker}
to count authorship and assignments as of November 19, 2026. Authored counts
include open and closed issues created by each member. Assigned counts include
issues that were closed and assigned to that member at the snapshot, including
issues opened by external reviewers. An issue with multiple assignees counts
for each assignee. Pull requests are counted in Section~6. The GitHub
accounts, in table order, are
\href{https://github.com/IbrahimT04}{IbrahimT04},
\href{https://github.com/chosterto}{chosterto},
\href{https://github.com/xSankalpaD}{xSankalpaD},
\href{https://github.com/vikramC04}{vikramC04}, and
\href{https://github.com/ArjunSasidaran}{ArjunSasidaran}.

Since the November 15 snapshot, Ibrahim opened 18 issues to track TA feedback
on the SRS and Hazard Analysis, and Vikram opened four issues covering the
Design Thinking slides, the backend module, PDF compilation, and module links.
The six earlier productivity report issues,
\href{https://github.com/SoftwareMandems/FPVR/issues/264}{\#264} through
\href{https://github.com/SoftwareMandems/FPVR/issues/269}{\#269}, are now
closed and contribute to Ibrahim's closed assignment count.

\section{CICD via GitHub Actions}

\wss{Say how CICD technology is used and will be used in your project.}

\wss{Provide links to your CICD yaml files}

\section{Custom Documents (CDs)}

\wss{What is the plan (as documented in TeamComposition.csv) for the team's
CDs?  Should the CDs be modified now that the team knows more about the
project?}

\section{Team Charter Trigger Items}

\wss{Provide a summary of the quantified triggers identified in the team's
charter.}

\wss{Provide a list of any violations of the triggers.  If the team wishes, the
violations can be summarized on aggregate, instead of naming specific team
members.}

\wss{Provide a plan to address the violations.  This could include revising the
triggers, if they are found to be too weak, strong or ambiguous.}

\section{Additional Productivity Metrics}

\wss{If your team has additional metrics of productivity, please feel free to
add them to this report.  If not, please explicitly state that there are no additional
metrics.}

\wss{Additional metrics can include things like code reviews done, pull requests
created, count of joining meetings late, count of number of times contributions
had to be corrected, number of internal deadlines missed, test cases written, etc.}

\wss{We are looking for data on these metrics, not just a list of additional
metrics the team is planning on tracking.  However, if all you have is a plan,
please share it here.}

\end{document}
